\section{The shift is a hyperbolic rotation}\label{sec:hyperbolic}

Theorem~\ref{thm:kernel} has an immediate consequence that identifies the shift
geometrically. Let $X$ be multiplication by $x$, so that $\cosh(\omega X)$ and
$\sinh(\omega X)$ map $C_c^\infty(I_L)$ to itself.

\begin{proposition}[Hyperbolic form of the shift]\label{prop:hyperbolic}
For $0\leq\omega\leq\frac12$ and $f\in C_c^\infty(I_L)$,
\begin{equation}\label{eq:hyperbolic}
 Q_{\omega,L}[f]=Q_{0,L}\big[\cosh(\omega X)f\big]
 -Q_{0,L}\big[\sinh(\omega X)f\big].
\end{equation}
\end{proposition}

\begin{proof}
Write $g_c=\cosh(\omega X)f$, $g_s=\sinh(\omega X)f$, $G_c=E_Lg_c$,
$G_s=E_Lg_s$, and evaluate \eqref{eq:target} on each, term by term. For the
norms, $|\!\cosh(\omega x)|^2-|\!\sinh(\omega x)|^2=1$ pointwise, so
$\norm{g_c}^2-\norm{g_s}^2=\norm f^2$; this disposes of the contact term and of
the $2\norm F^2$ inside \eqref{eq:gamma-energy}. For a translation, at any
$r$,
\begin{align*}
 \re\ip{G_c}{U_rG_c}-\re\ip{G_s}{U_rG_s}
 &=\re\int\big[\cosh(\omega x)\cosh(\omega(x-r))\\
 &\qquad\qquad-\sinh(\omega x)\sinh(\omega(x-r))\big]
   \overline{F(x)}F(x-r)\dd x\\
 &=\cosh(\omega r)\,\re\ip F{U_rF},
\end{align*}
by the addition formula $\cosh a\cosh b-\sinh a\sinh b=\cosh(a-b)$. Applying this
at $r$ inside the $\gammak$ integral and at $r=m\log p$ in the prime sum
reproduces the first and last lines of \eqref{eq:shifted-target}. For the pole
form, $P_L[g]=2\iint\cosh(\frac{x-y}2)\overline{g(x)}g(y)$, and the same addition
formula applied to $\cosh(\omega x)\cosh(\omega y)-\sinh(\omega x)\sinh(\omega y)
=\cosh(\omega(x-y))$ reproduces the third line. Each difference is computed
inside a convergent integral, so no regularization is needed. Summing gives
\eqref{eq:shifted-target}, which is $Q_{\omega,L}[f]$ by
Theorem~\ref{thm:kernel}.
\end{proof}

Equation \eqref{eq:hyperbolic} is (1.10) of \cite{Background}; the proof above
derives it from the generator rather than from a kernel expansion, which is what
makes the next remark available.

\begin{remark}[The shift is an angle, not a coupling]\label{rem:angle}
The pair $(\cosh(\omega X),\sinh(\omega X))$ is a hyperbolic rotation by angle
$\omega$ with generator $X$. Two consequences are worth stating.

First, the shift is a \emph{geometric} deformation of the input, not a change of
coefficients, so it is the kind of parameter a line can carry. In gauge theory
the natural object with a hyperbolic angle attached to a line is a \emph{cusp};
this is where we would look first among specific theories, and the differential
equation a cusped line's expectation satisfies in its angle has exactly the
shape the proposal asks for \cite{CHMS}.

Second, the parameter is not free. A cusp angle is a geometric datum of the
contour, so a mechanism built on it cannot be dialled to fit; that is the
property the program wants, since any mechanism with a tunable margin uniform in
$L$ proves something false \cite{Selection}. The obstruction to state plainly is
one of type: a cusp anomalous dimension is a number, while $Q_{\omega,L}$ is a
form on an infinite-dimensional input space, so a line carrying it must carry an
infinite-dimensional label. This is Condition~\ref{cond:unbounded} below.
\end{remark}

\begin{remark}[Positivity at positive shift is a different statement]
\label{rem:not-criterion}
Proposition~\ref{prop:hyperbolic} exhibits $Q_{\omega,L}$ as a \emph{difference}
of two central energies. Nonnegativity of each does not give nonnegativity of the
difference, so positivity of the shifted form at all shifts is strictly stronger
than the criterion \eqref{eq:criterion} and is not implied by the Riemann
hypothesis. The route of Proposition~\ref{prop:route} does not require it: what
it requires is $\norm{V_{\omega_j,L_j}}\leq1$, a statement about the cumulative
defect \eqref{eq:storage} rather than about its integrand.
\end{remark}
